\documentclass[11pt]{ctexart}
\usepackage[a4paper,margin=1in]{geometry}
\usepackage{graphicx}
\usepackage{xcolor}
\usepackage{hyperref}
\definecolor{refgray}{gray}{0.45}
\title{\textbf{Market Views｜全球投行叙事汇编}\\[0.4em]{\large 通胀降温主线确认，但地缘风险与AI资本开支回报验证制造短期张力}}
\author{KC桌面 自动生成}
\date{260720}
\begin{document}
\maketitle
\tableofcontents
\newpage
\section{一页摘要}
\begin{itemize}
  \item 全球宏观主线从通胀恐慌转向通胀确认降温，美国核心PCE被指高估真实通胀，为央行提供政策空间，但中东局势推升油价构成短期上行风险。
  \item 中国二季度GDP增速放缓至4.3\%，低于目标下限，三季度财政加速成为维持全年目标的必要选项，出口量价背离与内需疲弱裂口扩大。
  \item AI基础设施资本开支仍处上行周期，超大规模云厂商2026年支出预计接近8700亿美元，但市场对货币化路径的追问日益迫切，分歧集中在供给瓶颈与需求结构。
  \item 中国AI基础设施投资加速，GS上修2026-2028年云厂商资本开支预测37\%-55\%，驱动本土芯片出货量CAGR达92\%。
  \item 中国电池消费税重启加速行业洗牌，头部企业凭借利润垫与海外布局韧性凸显，二线厂商面临整合压力，成本增幅1-3\%不足以改变需求曲线。
  \item 铜市场实物基本面收紧，但宏观不确定性压制价格上行空间，AI叙事正成为铜价定价的新驱动力，与实物库存紧张形成双轨叙事。
  \item 中国消费动能环比放缓，但收入修复与结构性升级并存，茅台打破惯例一年内两次调价，市场对其定价权从供需驱动转向战略主导的解读存在分歧。
  \item 欧洲电气化目标与数据中心、热泵、电动车渗透率拐点共振，电网投资缺口2.2-3.5万亿欧元成为确定性增长主线，但电费上涨可能抑制终端进程。
\end{itemize}
\section{全球投行叙事汇编}
今日纳入 31 篇投行/券商研究，按 4 个市场、资产与产业主题系统整理。
\newpage
\subsection{宏观与策略：通胀降温主线与地缘扰动下的资产定价}
\textbf{全球宏观主线正从通胀恐慌转向通胀确认降温，但地缘政治风险（中东、贸易）与结构性分化（AI估值、中国内需）制造了短期扰动与长期机会之间的张力。市场共识是通胀趋势性回落为央行提供政策空间，分歧在于油价上行风险与AI调整的深度是否足以逆转这一进程。}
\subsubsection*{主流共识}
\begin{itemize}
  \item 全球除美国外核心通胀已接近目标，美国核心PCE高估真实通胀，通胀降温是中期主线。
  \item 中国二季度GDP增速放缓至4.3\%，低于目标下限，三季度财政加速成为维持全年目标的必要选项。
  \item AI相关板块经历深度调整（KOSPI -25\%，SOX -20\%），但盈利支撑依然强劲，基本面尚未逆转。
\end{itemize}
\subsubsection*{分歧与条件差异}
\begin{itemize}
  \item 油价风险：中东局势升级推高布伦特油价，但GS认为供需两端都存在快速逆转可能，通胀影响可控；NOM的TACO指标升至2.0标准差，但认为立即爆发概率低。
  \item 韩国财富效应：市场上涨带来GDP约17\%的财富增值，但财富集中于高储蓄老年群体，消费转化率仅1.6\%，且加息可能抵消这一微弱效应。
  \item 中国贸易数据：6月出口同比+27\%创纪录，但GS拆解显示量价背离（半导体进口量仅+6.6\%），内需疲弱（零售+1.0\%）与生产强劲（工业+5.3\%）的裂口扩大。
\end{itemize}
\subsubsection*{机构观点}
\paragraph{Bernstein} 印度市场估值已重置至18.6倍远期PE（低于十年均值0.7个标准差），金融和科技板块PE回到疫情低位。外资流出逆转（7月净流入18亿美元），大盘股盈利修正出现早期复苏信号。维持中性，但机会组合正在变宽。
{\small\color{refgray}数据：MSCI印度远期PE 18.6x，低于十年均值0.7个标准差\par}
{\small\color{refgray}数据：金融和科技板块PE低于均值1.9和1.5个标准差\par}
{\small\color{refgray}数据：7月外资净流入18亿美元，全球新兴市场基金低配从-2.7\%收窄至-2.1\%\par}
{\small\color{refgray}数据：长期盈利增长预期降至10.6\%，低于历史均值\par}
{\small\color{refgray}边际变化：印度估值从23.2倍PE高点回落至18.6倍，外资流出趋势逆转，大盘股盈利上修信号出现，机会从防御转向价值。\par}
\paragraph{GS} 韩国市场上涨为家庭带来GDP约17\%的财富增值，但财富集中于前20\%高收入家庭（持有64\%股票资产）和50岁以上群体（73\%），消费转化率仅1.6\%。高杠杆（偿债负担占收入27\%）和高储蓄率（60岁以上37\%）制约财富效应，加息可能完全抵消消费拉动。
{\small\color{refgray}数据：韩国市场上涨带来相当于GDP 17\%的家庭财富增值\par}
{\small\color{refgray}数据：每100韩元财富增值仅带动1.6韩元额外消费\par}
{\small\color{refgray}数据：前20\%高收入家庭持有64\%股票资产\par}
{\small\color{refgray}数据：50岁以上群体持有73\%家庭股票财富\par}
{\small\color{refgray}边际变化：市场上涨幅度创历史新高，但财富效应传导路径被结构性因素（集中度、高杠杆、高储蓄）显著削弱，消费拉动仅0.3个百分点。\par}
\paragraph{GS} 中国二季度GDP同比4.3\%，低于市场预期4.5\%和政府目标下限。生产端（工业+5.3\%）与需求端（零售+1.0\%，固投-5.7\%）裂口扩大。6月贸易顺差创纪录1256亿美元，但半导体进口量仅+6.6\%，量价背离。三季度财政加速（8000亿政策性金融工具）是维持全年目标的关键。
{\small\color{refgray}数据：二季度实际GDP同比4.3\%，低于市场预期4.5\%\par}
{\small\color{refgray}数据：6月工业增加值同比+5.3\%，零售销售+1.0\%\par}
{\small\color{refgray}数据：固定资产投资年初至今同比-5.7\%\par}
{\small\color{refgray}数据：6月贸易顺差1256亿美元，创历史新高\par}
{\small\color{refgray}边际变化：GDP增速从Q1的5.0\%放缓至Q2的4.3\%，低于政府目标下限，触发三季度财政加速预期；贸易数据名义强劲但量价背离，内需疲弱成为核心矛盾。\par}
\paragraph{GS} 中东局势升级推高布伦特油价至预测路径上方，但通胀降温才是主线。G10除美国外核心通胀仅2.1\%，美国核心PCE 3.3\%存在测量偏差。油价风险双向存在，供需两端均可快速逆转。美联储7月加息概率基本消失，下半年美国消费预计放缓。
{\small\color{refgray}数据：G10除美国外核心通胀2.1\%，接近目标\par}
{\small\color{refgray}数据：美国核心PCE 3.3\%，存在测量偏差\par}
{\small\color{refgray}数据：布伦特期货突破GS此前预测的2026Q4 80美元/桶\par}
{\small\color{refgray}数据：6月核心PCE预计环比有限，同比回落\par}
{\small\color{refgray}边际变化：油价风险偏向上行，但GS强调通胀趋势性降温才是中期主线；美国核心PCE高估真实通胀，全球其他地区通胀已接近目标。\par}
\paragraph{JPM} AI动量退潮进入成熟阶段，KOSPI -25\%，SOX -20\%，但MSCI全球指数仅较高点回落1-2\%。半导体价格与盈利出现罕见裂口，盈利相对表现坚挺，产能增量2028年前有限。通胀回落（美国CPI三个月年化从8.2\%降至2.8\%）改善利率环境，支持市场广度扩大。
{\small\color{refgray}数据：KOSPI较高点下跌25\%，SOX下跌20\%\par}
{\small\color{refgray}数据：MSCI全球指数仅较高点回落1-2\%\par}
{\small\color{refgray}数据：美国CPI三个月年化从5月8.2\%降至6月2.8\%\par}
{\small\color{refgray}数据：半导体RSI接近超卖区域\par}
{\small\color{refgray}边际变化：AI板块调整深度加大，但盈利支撑未变，通胀快速回落改善利率环境，JPM认为调整不会演变为持续下跌，半导体可能很快找到买盘。\par}
\paragraph{GS} 洛杉矶港进口未来两周先升后降（7月24日当周环比+3\%，7月31日当周环比-7\%），并非趋势反转而是高波动下的喘息窗口。中国出发货量同比-25.5\%，环比-19\%，降幅扩大。贸易政策不确定性从一次性冲击转向持续性规划难题。
{\small\color{refgray}数据：中国出发至美国集装箱船同比-25.5\%，环比-19\%\par}
{\small\color{refgray}数据：洛杉矶港7月24日当周TEU预计环比+3\%，7月31日当周-7\%\par}
{\small\color{refgray}数据：中国主要港口吞吐量环比-15\%，同比-11\%\par}
{\small\color{refgray}数据：中国/东亚至美西运价同比+204\%\par}
{\small\color{refgray}边际变化：洛杉矶港到港数据出现短暂反弹，但中国出发货量加速下滑，贸易政策不确定性从一次性冲击转向持续性规划难题。\par}
\paragraph{NOM} TACO不确定性指标从接近0标准差升至2.0标准差，主因美伊停火备忘录失效（6月17日签署，7月7日重新攻击）及霍尔木兹海峡事实性关闭。但立即爆发概率低，油气价格需再涨10\%才触及3.2标准差阈值。当前更似博弈升级而非全面冲突。
{\small\color{refgray}数据：TACO指标从7月6日约0sd升至2.0sd\par}
{\small\color{refgray}数据：油气价格需再涨10\%才触及3.2sd阈值\par}
{\small\color{refgray}数据：美伊停火备忘录6月17日签署，7月7日失效\par}
{\small\color{refgray}数据：霍尔木兹海峡事实性关闭\par}
{\small\color{refgray}边际变化：中东地缘不确定性快速攀升，但NOM认为尚未触发紧急应对机制，当前指标反映不确定性积累而非临界点。\par}
\paragraph{NOM} 美元/人民币中间价预测模型下调110基点至6.7824，计入逆周期因子后为6.7862，较前一日实际中间价低72基点。模型误差收敛，外部因素边际影响减弱，定价预期正在重新校准。
{\small\color{refgray}数据：模型预测值从6.7934下调至6.7824，幅度110基点\par}
{\small\color{refgray}数据：计入逆周期因子后预测6.7862，较实际中间价低72基点\par}
{\small\color{refgray}数据：模型误差收敛，定价预期重新校准\par}
{\small\color{refgray}边际变化：中间价预测下调，逆周期因子修正方向转负，模型误差收敛表明市场对官方定价意图的理解趋于一致。\par}
\subsubsection*{关键数据}
\begin{itemize}
  \item MSCI印度远期PE 18.6x，低于十年均值0.7个标准差
  \item 韩国市场上涨带来GDP 17\%财富增值，消费转化率仅1.6\%
  \item 中国二季度GDP 4.3\%，低于目标下限，6月贸易顺差1256亿美元创纪录
  \item G10除美国外核心通胀2.1\%，美国核心PCE 3.3\%
  \item KOSPI -25\%，SOX -20\%，MSCI全球指数仅回落1-2\%
  \item 中国出发至美国集装箱船同比-25.5\%
  \item TACO指标升至2.0sd，油气价格需再涨10\%至3.2sd
  \item 美元/人民币中间价预测下调110基点至6.7824
\end{itemize}
\begin{figure}[htbp]
\centering
\includegraphics[width=0.88\linewidth]{figures/fig_001.jpg}
\caption{Exhibit 1 - Exhibit 1-Exhibit 3) EXHIBIT 1: India valuations are now cheap both on absolute and relative basis, which trade at 18.6x on fwd PE and at 3.3x on PB. Relative to EM, India's PB has fallen near COVID lows India relative to EM-12m}
\end{figure}
\begin{figure}[htbp]
\centering
\includegraphics[width=0.88\linewidth]{figures/fig_027.jpg}
\caption{Exhibit 1 - Exhibit 2: Korean households' equity assets could have risen by \$15\textbackslash{}\%\$ of GDP from the end of last year Percent of GDP}
\end{figure}
\begin{figure}[htbp]
\centering
\includegraphics[width=0.88\linewidth]{figures/fig_062.jpg}
\caption{Figure 1 - Figure 2: SOX Index ytd}
\end{figure}
\begin{figure}[htbp]
\centering
\includegraphics[width=0.88\linewidth]{figures/fig_067.jpg}
\caption{Exhibit 3 - Exhibit 4: TEU growth was decelerated in the most recent week Daily TEU from China to the US, YoY Data as of 7/16/2026. Represents the aggregated container volume, measured in twenty-foot equivalent units (TEU), of vessels depart}
\end{figure}
{\small\color{refgray}本节综合 8 篇研究；涉及 Bernstein、GS、JPM、NOM。}
\newpage
\subsection{AI基础设施：资本开支加速与回报验证的十字路口}
\textbf{AI基础设施投资正从‘要不要投’进入‘投了能赚回多少’的验证阶段，超大规模云厂商资本开支指引上修与本土芯片替代加速构成核心主线，但市场对货币化路径的追问日益迫切，分歧集中在供给瓶颈（存储器、光模块）与需求结构（K型消费、AI推理成本）的边际变化上。}
\subsubsection*{主流共识}
\begin{itemize}
  \item AI资本开支仍处于上行周期，超大规模云厂商2026年合计支出预计接近8700亿美元，同比增长约77\%，其中超大规模厂商贡献约7500亿美元。
  \item 中国AI基础设施投资加速，GS将2026-2028年中国主要云厂商AI资本开支预测分别上修37\%、44\%和55\%，驱动本土芯片出货量2025-2030年CAGR达92\%。
  \item 800G/1.6T光模块出货爬坡是AI硬件放量的关键瓶颈之一，光学芯片供应改善推动第二季度净利润指引超预期44\%。
\end{itemize}
\subsubsection*{分歧与条件差异}
\begin{itemize}
  \item 存储器周期性质：MS认为当前短缺是结构性供给瓶颈（数据中心内存价格三季度环比涨超25\%），而非周期性见顶；市场部分观点担忧长期协议和去规格化压低价格峰值。
  \item AI推理成本与货币化时间差：Bernstein认为token消耗量爆发需代理交易量起飞，市场对AI成本担忧过度；JPM则强调市场对货币化追问正变得更具体和迫切。
  \item 人形机器人技术路线：Bernstein主张从VLA向世界动作模型（WAM）迁移将重塑竞争壁垒，削弱数据垄断优势；但行业仍处于早期，长序列自主规划与任务泛化尚未突破。
  \item 在线杂货增长驱动力：GS认为增长主力来自现有用户购买频率提升（近50\%消费者增加快速配送订单），而非新增用户；市场部分观点仍聚焦用户获取竞争。
\end{itemize}
\subsubsection*{机构观点}
\paragraph{Bernstein} 人形机器人行业正从VLA模型向世界动作模型（WAM）迁移，核心差异在于WAM先预测物理可行结果再生成动作序列，类似人类行动前在脑中模拟后果。这一转变将削弱‘数据垄断’壁垒，模型架构和数据整合能力成为新护城河。当前行业在运动控制（Unitree机器人能追野猪）、操作（连续数小时自主分拣）上取得进展，但长序列自主规划和任务泛化仍是下一阶段战场。
{\small\color{refgray}数据：WAM分为渲染型、模拟型、规划型三类\par}
{\small\color{refgray}数据：Physical Intelligence的π0.7模型引入世界模型后复杂任务表现提升\par}
{\small\color{refgray}数据：两年前平地稳定行走还是成就，如今Unitree机器人能在波兰森林追逐野猪\par}
{\small\color{refgray}边际变化：行业焦点从‘能不能动’转向‘能不能想’，VLA到WAM的算法升级可能重新定义竞争逻辑\par}
\paragraph{Bernstein} 中国互联网板块估值回调至2022-2023年低谷水平后，不确定性回报比明显改善。AI推理成本担忧过度：聊天式对话仅消耗数百token，而代理式交易需数万个token，token消耗量爆发需代理交易量起飞。腾讯Hy3模型发布、阿里云收入增速加速至45\%左右、电商利润好于预期，推动近期反弹。未来几个月战术性环境建设性，微信智能助手Xiaowei逐步开放，AI变现路径渐清晰。
{\small\color{refgray}数据：覆盖公司年初至今平均下跌约16\%\par}
{\small\color{refgray}数据：阿里云6月季度收入增速约45\%，增量现金利润率35-40\%\par}
{\small\color{refgray}数据：腾讯估值倍数约11-12倍\par}
{\small\color{refgray}数据：阿里云下半年计算价格将继续上调\par}
{\small\color{refgray}边际变化：市场从‘坏消息充分定价’转向关注AI变现路径，微信智能助手和阿里云提价成为新催化剂\par}
\paragraph{GS} 美国在线杂货配送行业增长加速，2025-2030年CAGR达11\%，2030年市场规模接近4000亿美元。当日配送增速最快（+16\% CAGR），核心驱动力来自现有用户购买频率提升而非新增用户——近50\%消费者增加快速配送订单次数。四家头部公司（亚马逊、DoorDash、Uber、Instacart）在足够大的空白市场中各自发展，竞争未挤压增长。配送时效从‘小时级’向‘分钟级’跨越是频率提升的物理前提。
{\small\color{refgray}数据：2025-2030年行业CAGR +11\%\par}
{\small\color{refgray}数据：当日配送CAGR +16\%，远超自提（+4\%）和邮寄到家（+9\%）\par}
{\small\color{refgray}数据：近50\%消费者增加快速配送订单次数\par}
{\small\color{refgray}数据：2030年市场规模接近4000亿美元\par}
{\small\color{refgray}边际变化：增长驱动力从‘拉新’转向‘复购’，配送时效缩短和自动化履约中心规模化是频率提升前提\par}
\paragraph{GS} 全球夏季出行需求稳定健康，6月预订趋势超季初预期，但结构呈K型分化：高端市场定价上移，大众消费者削减外出就餐、网购等开支‘挤出’旅行预算。单次旅行支出上升，预订窗口两极分化（远期与临期并存）。AI在旅行规划中使用率显著提升（1H26美国超56\%休闲旅客使用AI规划旅行，去年同期33\%），但通过AI完成预订比例仍低，品牌信任是前提。
{\small\color{refgray}数据：1H26美国超56\%休闲旅客使用AI规划旅行（去年同期33\%）\par}
{\small\color{refgray}数据：约三分之一受访者削减日常开支为旅行腾预算\par}
{\small\color{refgray}数据：北美主办城市赛事期间短租均价同比上涨60\%以上\par}
{\small\color{refgray}数据：部分城市入住率达90\%\par}
{\small\color{refgray}边际变化：消费者预算再分配而非消费升级，AI规划使用率翻倍但预订转化仍低\par}
\paragraph{GS} 中国数据中心电气公司股价回调约29\%（同期沪深300跌9\%），核心担忧是技术路线切换和国内价格竞争，但基本面未变。海外扩张是估值修复关键：Kstar大功率UPS产线接近满负荷，获台湾客户集成方案订单、日本/欧洲客户认证及欧洲数据中心超2亿元订单；Megmeet的Vera Rubin机架电源出货加速，800VDC产品已送样。国内资本开支环比改善，但竞争激烈价格承压。
{\small\color{refgray}数据：Kstar大功率UPS产线接近满负荷运转\par}
{\small\color{refgray}数据：Kstar获欧洲数据中心超2亿元UPS订单\par}
{\small\color{refgray}数据：Kstar已推出AC-800VDC模块，三季度送样测试\par}
{\small\color{refgray}数据：Kehua获阿里超1亿元模块化配电+CDU订单\par}
{\small\color{refgray}边际变化：海外订单落地进度成为估值修复催化剂，800VDC产品商业化节奏决定中期增长空间\par}
\paragraph{GS} 一家中国光模块公司2026年第二季度净利润指引中值较预期高出44\%，核心驱动是800G/1.6T光模块出货加速，而非总量突然放大。光学芯片供应改善是爬坡前提，公司产能扩张配合节奏。2026年盈利预测上调7\%，2027-2028年仅上调2\%，说明近期供应改善对远期格局影响有限。毛利率从2023年31\%提升至2025年47.8\%，预计2026-2028年稳定在48\%以上。
{\small\color{refgray}数据：2Q26净利润中值较预期高出44\%\par}
{\small\color{refgray}数据：2Q26净利润同比增长78\%-163\%，环比增长68\%\par}
{\small\color{refgray}数据：2026年营收预期上调6\%至514.6亿\par}
{\small\color{refgray}数据：2026年净利润预期上调7\%至206.8亿\par}
{\small\color{refgray}边际变化：800G/1.6T出货加速驱动环比大幅改善，但长期增长逻辑将从规模扩张转向规模互联\par}
\paragraph{GS} 中国AI基础设施投资加速上行，本土芯片供应商份额提升刚刚开始。GS将2026-2028年中国主要云厂商AI资本开支预测分别上修37\%、44\%和55\%。一家本土GPU芯片公司收入预计2025-2030年CAGR达108\%，芯片出货量从2025年4.5万片攀升至2030年超100万片。全球领先AI芯片供应商在中国仍占超50\%市场份额，本土替代空间显著。
{\small\color{refgray}数据：2026-2028年中国主要云厂商AI资本开支预测上修37\%/44\%/55\%\par}
{\small\color{refgray}数据：本土GPU芯片公司收入2025-2030年CAGR 108\%\par}
{\small\color{refgray}数据：芯片出货量从2025年4.5万片增至2030年超100万片\par}
{\small\color{refgray}数据：全球领先AI芯片供应商在中国仍占超50\%市场份额\par}
{\small\color{refgray}边际变化：中国云厂商AI投资意愿强于市场预期，本土芯片替代从‘开始’进入‘加速’阶段\par}
\paragraph{JPM} AI基础设施叙事未变，但市场关注点从资本开支规模转向货币化回报。到2026年底AI相关资本开支预计接近8700亿美元（同比增长约77\%），超大规模厂商贡献约7500亿美元。JPM对2027年预测更积极：谷歌资本开支增长54\%至约3000亿美元，亚马逊和Meta分别增长42\%至3000亿和2000亿美元。动量因子从极端拥挤回落，但仓位尚未完全出清，波动可能持续。
{\small\color{refgray}数据：2026年底AI相关资本开支预计接近8700亿美元，同比增长约77\%\par}
{\small\color{refgray}数据：超大规模厂商贡献约7500亿美元\par}
{\small\color{refgray}数据：2027年谷歌资本开支预计增长54\%至约3000亿美元\par}
{\small\color{refgray}数据：2027年亚马逊和Meta资本开支分别增长42\%至3000亿和2000亿美元\par}
{\small\color{refgray}边际变化：市场焦点从‘投多少’转向‘赚回多少’，货币化信号若浮现将支撑盈利和自由现金流曲线\par}
\paragraph{MS} 存储器从AI需求受益者变为AI供给瓶颈。数据中心内存价格三季度环比上涨至少25\%，高于预期，短缺没有缓解反而加剧。AI消耗DRAM过多，PC和智能手机出货受抑制。长期协议和去规格化不会终结周期方向，而是拉长周期形状——客户愿意为六周加急订单支付高于合同价的价格，反映真实产能挤兑而非囤货。
{\small\color{refgray}数据：数据中心内存价格三季度环比上涨至少25\%\par}
{\small\color{refgray}数据：客户愿为六周加急订单支付高于二季度合同价的价格\par}
{\small\color{refgray}数据：AI消耗DRAM过多，抑制PC和智能手机出货\par}
{\small\color{refgray}数据：英伟达削减机架中LPDDR5内存含量以优化现有约束\par}
{\small\color{refgray}边际变化：存储器从AI需求受益者变为供给瓶颈，短缺结构性而非周期性，长期协议拉长盈利可见度\par}
\subsubsection*{关键数据}
\begin{itemize}
  \item 2026年底AI相关资本开支预计接近8700亿美元，同比增长约77\%
  \item GS将2026-2028年中国主要云厂商AI资本开支预测上修37\%/44\%/55\%
  \item 一家中国光模块公司2Q26净利润中值较预期高出44\%，环比增长68\%
  \item 数据中心内存价格三季度环比上涨至少25\%
  \item 本土GPU芯片出货量从2025年4.5万片增至2030年超100万片，CAGR 92\%
  \item 美国在线杂货2025-2030年CAGR 11\%，当日配送CAGR 16\%
  \item 1H26美国超56\%休闲旅客使用AI规划旅行（去年同期33\%）
  \item Kstar获欧洲数据中心超2亿元UPS订单，大功率UPS产线接近满负荷
\end{itemize}
\begin{figure}[htbp]
\centering
\includegraphics[width=0.88\linewidth]{figures/fig_017.jpg}
\caption{Exhibit 2 - Exhibit 2: We expect the industry to reach \textbackslash{}\textasciitilde{}\textbackslash{}\$400bn in sales by 2030... US grocery spend (\textbackslash{}\$bn) \#\# Across the three key categories of online grocery spend, we forecast:}
\end{figure}
\begin{figure}[htbp]
\centering
\includegraphics[width=0.88\linewidth]{figures/fig_051.jpg}
\caption{Exhibit 1 - Exhibit 2: Earnings revision}
\end{figure}
\begin{figure}[htbp]
\centering
\includegraphics[width=0.88\linewidth]{figures/fig_072.jpg}
\caption{Figure 2 - Figure 2: AI Capex (LTM Actuals vs Estimates) Figure 3: Hyperscaler Income Statement}
\end{figure}
\begin{figure}[htbp]
\centering
\includegraphics[width=0.88\linewidth]{figures/fig_076.jpg}
\caption{Exhibit 1 - Exhibit 1: DRAM multiples have corrected sharply, spot pricing is continuing to move higher Aren't the long-term agreements capping upside? Our view is that the long-term agreements are important to the stocks, but more because t}
\end{figure}
{\small\color{refgray}本节综合 9 篇研究；涉及 Bernstein、GS、JPM、MS。}
\newpage
\subsection{能源与工业：电池税改重塑竞争格局，铜价叙事转向AI与实物双轨}
\textbf{中国电池消费税重启加速行业洗牌，头部企业凭借利润垫与海外布局韧性凸显，二线厂商面临整合压力；铜市场实物收紧与宏观情绪退潮并存，AI叙事正成为铜价定价的新驱动力；欧洲电气化计划开启电网投资长周期，储能与电力设备需求结构性跃升。}
\subsubsection*{主流共识}
\begin{itemize}
  \item 中国电池消费税（2\%-4\%）对终端需求影响有限，但将加速行业集中，头部企业（如CATL）凭借高单位利润和海外收入占比可有效消化成本冲击。
  \item 铜市场短期实物基本面收紧（中国库存降至16.7万吨、LME注销仓单上升），但宏观不确定性（中东局势、美联储政策）压制价格上行空间。
  \item 欧洲电气化目标（2040年电力占比46\%）与数据中心、热泵、电动车渗透率拐点共振，电网投资缺口（2.2-3.5万亿欧元）成为确定性增长主线。
\end{itemize}
\subsubsection*{分歧与条件差异}
\begin{itemize}
  \item 电池消费税影响程度：Bernstein认为成本增幅1-3\%不足以改变需求曲线，GS强调单位利润差异将导致企业间影响分化（CATL影响1-6\% vs 二线厂商8-82\%）。
  \item 铜价驱动因素权重：GS模型显示AI叙事解释力持续上升，但实物库存紧张主要由废铜替代而非终端需求拉动，市场对定价主导权存在分歧。
  \item 欧洲电力需求增速：GS预计2029年后年均4-5\%增长，但部分观点担忧电费上涨可能抑制终端电气化进程，实际增速存在不确定性。
\end{itemize}
\subsubsection*{机构观点}
\paragraph{Bernstein} 储能电池正从汽车供应链附属角色独立，三星SDI为特斯拉建专用储能产线（200Ah LFP 2026年底量产、300Ah 2027年Q2投产），LG供应谷歌最大光储项目，AI数据中心成为储能需求新增长极。中国电池消费税（2\%-4\%）对终端价格影响仅1-3\%，政策意图在于遏制产能过剩并引导向钠离子/固态电池切换，预计2030年中国10-20\%电池产量为固态或钠离子。瑞萨存储器接口调查影响可控，该业务仅占营收6-8\%，毛利率65\%由IP壁垒等结构性因素支撑，最差情景下对利润影响仅1.8\%。
{\small\color{refgray}数据：三星SDI储能产线投资3-5万亿韩元\par}
{\small\color{refgray}数据：LFP电芯价格79.3美元/kWh，NMC 101.5美元/kWh\par}
{\small\color{refgray}数据：瑞萨存储器接口业务毛利率65\%\par}
{\small\color{refgray}数据：最差情景对瑞萨利润影响1.8\%\par}
{\small\color{refgray}边际变化：储能电池从借用动力电池产能转向自建专用产线，消费税政策从全面免税转向技术差异化征税，存储器接口调查未改变结构性高利润基础。\par}
\paragraph{GS} 铜实物基本面收紧（中国库存16.7万吨、LME注销仓单上升），但宏观情绪退潮，AI叙事正成为铜价波动日益重要的驱动力（数据中心仅占全球铜需求1\%，但重塑了市场对未来电力需求的预期）。中国电池消费税下CATL最具韧性，2025年单位净利润109元/kWh，海外收入占比超30\%，4\%税率下盈利影响仅1-6\%；二线厂商（中创新航、瑞浦兰钧）在50\%转嫁假设下影响8-41\%。欧洲电气化计划（2040年电力占比46\%）需2.2-3.5万亿欧元投资，电网平均服役40年，2029年后电力需求年增4-5\%。
{\small\color{refgray}数据：中国铜库存16.7万吨，处于季节性低位\par}
{\small\color{refgray}数据：CATL单位净利润109元/kWh\par}
{\small\color{refgray}数据：欧洲未来十年电气化投资2.2-3.5万亿欧元\par}
{\small\color{refgray}数据：欧洲电力需求2029年后年增4-5\%\par}
{\small\color{refgray}边际变化：铜定价从'中国需求+美元'传统框架转向'AI叙事+实物库存'双轨；电池消费税放大企业单位利润差距，加速行业集中；欧洲电网从存量维护切换至增量建设。\par}
\paragraph{NOM} 中国电池消费税（2\%-4\%）对单车成本影响约400-500元（60kWh电池包），对整车制造影响有限。头部企业（如CATL）可通过成本转嫁或优化消化冲击，二线厂商议价能力弱、利润压力更大。出口导向型厂商及钠离子/固态电池布局者受益于免税政策。政策节奏符合预期，是税收优惠退坡的延续。
{\small\color{refgray}数据：单车成本增加400-500元\par}
{\small\color{refgray}数据：电芯价格400元/kWh，2\%税率对应7元/kWh税负\par}
{\small\color{refgray}边际变化：电池行业从政策红利驱动转向市场化竞争，二线厂商利润空间被进一步压缩，出口与新技术路线获得差异化优势。\par}
\subsubsection*{关键数据}
\begin{itemize}
  \item 中国电池消费税税率2\%（2026年9月）升至4\%（2027年9月）
  \item CATL单位净利润109元/kWh vs 二线厂商8元/kWh
  \item 中国铜库存16.7万吨，季节性低位
  \item 欧洲电网平均服役40年，未来十年投资2.2-3.5万亿欧元
  \item 欧洲电力需求2029年后年增4-5\%
  \item 三星SDI储能产线投资3-5万亿韩元
  \item 瑞萨存储器接口业务毛利率65\%，最差情景利润影响1.8\%
  \item 数据中心仅占全球铜需求1\%，但AI叙事重塑定价
\end{itemize}
\begin{figure}[htbp]
\centering
\includegraphics[width=0.88\linewidth]{figures/fig_050.jpg}
\caption{Exhibit 1 - Exhibit 1: A 4\% tax rate would imply an estimated unit tax of Rmb11–24/kWh for our covered companies, compared with their 2025 unit net profit of Rmb8–109/kWh.}
\end{figure}
\begin{figure}[htbp]
\centering
\includegraphics[width=0.88\linewidth]{figures/fig_037.jpg}
\caption{Exhibit 1 - Exhibit 3: China Copper Inventories Have Fallen to Seasonal Lows}
\end{figure}
\begin{figure}[htbp]
\centering
\includegraphics[width=0.88\linewidth]{figures/fig_016.jpg}
\caption{EXHIBIT 1 - EXHIBIT 1: Scenario Analysis on Renesas' Memory interface Business}
\end{figure}
\begin{figure}[htbp]
\centering
\includegraphics[width=0.88\linewidth]{figures/fig_038.jpg}
\caption{Exhibit 1 - Exhibit 3: China Copper Inventories Have Fallen to Seasonal Lows Exhibit 4: Scrap Tightness Has Shifted Fabricator Demand Toward Refined Cathode}
\end{figure}
{\small\color{refgray}本节综合 7 篇研究；涉及 Bernstein、GS、NOM。}
\newpage
\subsection{消费与金融：结构分化与主动调整}
\textbf{中国消费动能环比放缓，但收入修复与结构性升级并存；金融板块中，区域银行主动优化资产负债表，茅台调价频率加快，黄金珠宝短期承压但品牌升级持续。}
\subsubsection*{主流共识}
\begin{itemize}
  \item 中国消费总量增速放缓，但居民收入温和修复，储蓄倾向仍高，消费意愿谨慎。
  \item 啤酒行业8-10元价格带仍是高端化核心战场，大单品策略持续有效。
  \item 运动服饰行业整体增速放缓，但多品牌组合和运营调整可带来相对韧性。
\end{itemize}
\subsubsection*{分歧与条件差异}
\begin{itemize}
  \item 茅台一年内两次调价，打破历史惯例，市场对其定价权从供需驱动转向战略主导的解读存在分歧。
  \item 金价波动对黄金珠宝销售的影响是短期扰动还是长期趋势，机构看法不一。
  \item 区域银行主动压缩消费贷款、转向商业贷款的策略，市场对其盈利修复节奏的预期存在差异。
\end{itemize}
\subsubsection*{机构观点}
\paragraph{GS} 美国区域银行二季度财报显示，TFC、FITB、RF的PPNR均略超预期，但市场对指引细节更敏感。TFC下调2026年NII预期至持平/微降，同时上调手续费收入增速至10\%，并主动压缩低收益消费贷款，将资本转向商业贷款和交易业务，ROTCE预期从14\%上调至14\%以上。FITB并购CMA协同效应有望超预期，消费者存款增长\$25亿（vs预期\$10亿），但Q3 PPNR指引偏弱。RF手续费增长指引调至区间下沿。GS认为，资产结构优化正在对冲利率不确定性，估值修复空间存在。
{\small\color{refgray}数据：TFC 2026年NII预期从温和增长下调至持平/微降\par}
{\small\color{refgray}数据：TFC手续费收入增速从高个位数上调至10\%\par}
{\small\color{refgray}数据：TFC 2026年ROTCE预期从14\%上调至14\%以上\par}
{\small\color{refgray}数据：FITB消费者存款增长\$25亿，超预期\$10亿\par}
{\small\color{refgray}边际变化：市场此前关注利率环境改善对NII的提振，但GS指出银行正主动调整资产结构，用非利息收入弥补NII缺口，并上调ROTCE指引，显示盈利修复路径从被动等待转向主动优化。\par}
\paragraph{GS} 中国消费二季度出现收入改善但支出观望的分化：居民可支配收入名义同比增速从4.9\%升至5.6\%，社零同比增速却从2.4\%降至0.2\%。GS工资追踪指标显示城镇居民工资同比增速从4.3\%升至4.7\%，但居民储蓄率从32.3\%微升至32.6\%，超额存款达61万亿元。消费结构上，食品支出放缓，教育文化娱乐回升；以旧换新政策对家电仍有拉动，但对汽车提振减弱，二季度汽车销量低于去年同期。
{\small\color{refgray}数据：居民可支配收入名义同比增速Q1 4.9\%→Q2 5.6\%\par}
{\small\color{refgray}数据：社零同比增速Q1 2.4\%→Q2 0.2\%\par}
{\small\color{refgray}数据：城镇居民工资同比增速Q1 4.3\%→Q2 4.7\%\par}
{\small\color{refgray}数据：居民储蓄率32.3\%→32.6\%，超额存款61万亿元\par}
{\small\color{refgray}边际变化：此前市场关注收入修复能否带动消费反弹，但GS指出收入改善并未转化为支出增长，储蓄倾向仍高，消费动能环比放缓，结构上服务消费优于商品消费。\par}
\paragraph{GS} 中国啤酒行业结构性分化加剧，8-10元价格带仍是高端化核心战场。燕京U8 2025年销量达90万千升，2026年Q1仍维持近30\%同比增长，7月后势头加速，持续从青岛经典、雪花清爽等竞品夺取份额。华润啤酒同期销量可实现低个位数增长，而青岛啤酒和重庆啤酒分别下滑约4\%，百威中国录得约10\%降幅。大单品策略在总量见顶市场中仍有效，集中资源打透一个价格带比分散布局更高效。
{\small\color{refgray}数据：燕京U8 2025年销量约90万千升，2019年上市以来CAGR约55\%\par}
{\small\color{refgray}数据：2026年Q1 U8同比增长近30\%\par}
{\small\color{refgray}数据：华润啤酒同期销量低个位数增长\par}
{\small\color{refgray}数据：青岛啤酒和重庆啤酒分别下滑约4\%\par}
{\small\color{refgray}边际变化：此前市场认为啤酒高端化整体放缓，但GS指出8-10元价格带的大单品仍在高速渗透，头部品牌通过集中资源实现份额扩张，与行业总量疲软形成对比。\par}
\paragraph{GS} 贵州茅台2026年7月完成年内第二次飞天调价，出厂价累计上调17\%，建议零售价累计上调9\%，为历史首次一年内两轮调价。7月调价覆盖批发、直销、团购、电商全渠道，团购渠道建议零售价升至1,719元/瓶。GS估算，按渠道比例分摊后，2026年营收增量约1.2\%，利润增量约1.6\%；2027年贡献扩大至约3\%营收和4\%以上利润。调价后原箱飞天批价从1,650元升至1,720元，散瓶从1,630元升至1,680元。
{\small\color{refgray}数据：飞天出厂价累计上调17\%（1,269→1,369元/瓶）\par}
{\small\color{refgray}数据：建议零售价累计上调9\%（1,539→1,639元/瓶）\par}
{\small\color{refgray}数据：团购渠道建议零售价升至1,719元/瓶\par}
{\small\color{refgray}数据：2026年营收增量约1.2\%，利润增量约1.6\%\par}
{\small\color{refgray}边际变化：茅台打破过去调价需间隔数年且依赖行业上行周期的惯例，在库存调整未结束时主动提价，显示定价权从市场供需转向公司战略主导，调价频率本身成为关键信号。\par}
\paragraph{NOM} 安踏体育二季度运营数据显示，安踏核心品牌和FILA均录得低个位数同比增长，其他品牌（迪桑特、可隆等）增长25-30\%，跑赢运动服饰板块。上半年整体安踏核心/FILA/其他品牌分别增长中个位数/中个位数/35-40\%。公司主动进行运营调整：更换安踏主品牌CEO、推出安踏SV/超级安踏新店型、开设安踏Market大店、狼爪新店铺开，主品牌线上销售已恢复双位数增长。
{\small\color{refgray}数据：安踏核心品牌和FILA二季度低个位数同比增长\par}
{\small\color{refgray}数据：其他品牌二季度增长25-30\%\par}
{\small\color{refgray}数据：上半年其他品牌增长35-40\%\par}
{\small\color{refgray}数据：主品牌线上销售二季度双位数增长\par}
{\small\color{refgray}边际变化：此前市场担忧运动服饰行业整体减速，但NOM指出安踏通过多品牌策略和主动运营调整（换帅、新店型、线上恢复）展现出超出平均水平的韧性，横向跑赢板块。\par}
\paragraph{NOM} 老铺黄金二季度销售受金价波动影响，但公司已采取促销、新品、高端客户定向支持等措施，效果初步显现。门店升级持续推进：深圳万象天地、上海国金中心、上海恒隆广场等搬迁或扩建，增加服务面积和差异化体验。NOM调整FY26-28F收入/盈利预测5-6\%/9-11\%，但认为若金价趋于稳定，观望情绪可能缓解，门店网络和客户体验能力将成为长期竞争优势。
{\small\color{refgray}数据：二季度金价波动直接影响销售\par}
{\small\color{refgray}数据：FY26-28F收入预测下调5-6\%\par}
{\small\color{refgray}数据：FY26-28F盈利预测下调9-11\%\par}
{\small\color{refgray}数据：目标价从HKD 1,114调整至HKD 905\par}
{\small\color{refgray}边际变化：此前市场关注金价上涨对黄金珠宝的利好，但NOM指出金价节奏变化（波动而非单边上涨）反而抑制销售，公司应对措施从被动等待转向主动促销和门店升级，长期品牌定位不变。\par}
\paragraph{GS} 澳大利亚金矿行业正在经历一次重要的资产整合。GS在最新报告中重新评估了Genesis Minerals与Vault Minerals的合并方案，并基于其自下而上的矿山、选厂和成本模型，给出了一个核心判断：这笔交易在财务上具有增值潜力，但真正的价值来源并非简单的规模扩张，而是资本配置路径的转变——通过避免新建选厂来降低执行不确定性，同时释放现有设施的产能弹性。
{\small\color{refgray}数据：GS认为，合并后的实体在近两年内可将黄金产量提升至约69万至78万盎司，全维持成本降至约每盎司2850澳元，较合并前的独立预测低约10\%。同期息税折旧摊销前利润有望提升15\%至25\%。这些数字背后，是报告试图回答的一个更根本的问题：在金价波动和项目执行环境日益复杂的背景下，什么样的增长策略才是可持续的。\par}
{\small\color{refgray}数据：GS的分析显示，合并后产量提升的主要驱动力并非新矿山的开发，而是对现有加工能力的重新调度。具体而言，Tower Hill矿区的矿石将被运往King of the Hills选厂处理，从而避免新建一座选厂。GS估算，仅此一项即可节省约3.5亿澳元的资本支出，并规避当前西澳项目执行环境中的工期和成本不确定性。\par}
\subsubsection*{关键数据}
\begin{itemize}
  \item 居民可支配收入名义同比增速Q1 4.9\%→Q2 5.6\%，社零同比增速Q1 2.4\%→Q2 0.2\%
  \item 燕京U8 2025年销量90万千升，2026年Q1同比增长近30\%
  \item 飞天茅台出厂价累计上调17\%，2026年利润增量约1.6\%
  \item 安踏其他品牌二季度增长25-30\%，主品牌线上恢复双位数增长
  \item 老铺黄金FY26-28F盈利预测下调9-11\%，门店持续升级
  \item TFC 2026年ROTCE预期从14\%上调至14\%以上，手续费收入增速上调至10\%
  \item 居民储蓄率32.3\%→32.6\%，超额存款61万亿元
\end{itemize}
\begin{figure}[htbp]
\centering
\includegraphics[width=0.88\linewidth]{figures/fig_042.jpg}
\caption{Exhibit 1 - Exhibit 1: Premium beer SKU map Exhibit 2: We expect premiumization to remain intact in 2025-28E Exhibit 3: Yanjing significantly gained share in the premium segment over the past five years on U8 success}
\end{figure}
\begin{figure}[htbp]
\centering
\includegraphics[width=0.88\linewidth]{figures/fig_045.jpg}
\caption{Exhibit 1 - Exhibit 1: Year-over-year growth in household disposable income per capita rose in Q2 2026 Exhibit 2: Year-over-year growth in household consumption per capita increased slightly in Q2 2026 Exhibit 3: The boost from the consume}
\end{figure}
\begin{figure}[htbp]
\centering
\includegraphics[width=0.88\linewidth]{figures/fig_047.jpg}
\caption{Exhibit 2 - Exhibit 5: Retail sales growth decelerated in Q2 mainly on weaker goods sales}
\end{figure}
\begin{figure}[htbp]
\centering
\includegraphics[width=0.88\linewidth]{figures/fig_032.jpg}
\caption{Exhibit 2 - Exhibit 3: ...with the Leonora Hub adding 50-100kozpa to \textbackslash{}\textasciitilde{}430-500kozpa Exhibit 4: Increased production and reduced milling costs (with no Tower Hill mill) see costs fall near-term...}
\end{figure}
{\small\color{refgray}本节综合 7 篇研究；涉及 GS、NOM。}
\newpage
\section{辅助信号｜战略咨询}
本板块整合波士顿咨询三份报告，聚焦AI在政府信任修复、国防航空维护及欧洲军工制造瓶颈中的应用与挑战。
\subsection{AI重塑公共服务信任与效率}
\textbf{公民对政府使用生成式AI的接受度高于预期，但信任修复需从低风险场景切入，并优先解决监管框架与AI素养问题。}
\begin{itemize}
  \item 全球数字政府服务净满意度仅微增2个百分点至65\%，但75\%的公民期望政府服务达到顶级私企水平，满意度提升速度远低于预期。
  \item 公民对GenAI的抽象焦虑（51\%认为AI发展太快）与具体场景接受度存在分界线：63\%愿意通过GenAI获取简单服务，71\%接受多语言沟通，69\%支持政府雇员使用AI辅助工具。
  \item 在“用GenAI自动化服务准入决策”和“监控公众情绪”等高不确定性场景上，各司法管辖区接受度差异显著，政府需根据本地社会许可边界设计用例。
  \item 报告提出三步走策略：先通过小规模、高透明度实践积累信任，再用制度和规则锁定信任，最后才有资格讨论大规模创新，顺序不能颠倒。
\end{itemize}
\begin{figure}[htbp]
\centering
\includegraphics[width=0.88\linewidth]{figures/fig_106.jpg}
\caption{Exhibit 1 - Exhibit 1: We surveyed 41,600 regular internet users across 48 jurisdictions \$\textasciicircum{}\{1\}\$ Special Administration Region of the People's Republic of China In conducting this research, we also held interviews and gathered qualitative inp}
\end{figure}
\subsection{AI驱动复杂装备维护与供应链效率提升}
\textbf{全球空军战机任务出勤率长期低于70\%，问题不在投入规模而在流程效率，AI可通过知识整合、备件可视化和流程优化重构维护体系。}
\begin{itemize}
  \item 美、法、德战机任务可用率常规性低于70\%，英国2024年F-35机队仅三分之一具备完全任务能力，合格维护人员不足是核心硬约束。
  \item GenAI可将分散的纸质手册和技术文档整合为交互式界面，使经验不足的维护人员快速定位故障并列出所需工具和备件，将“十年经验”压缩为一次查询。
  \item 备件短缺并非主因，备件“在某个角落但找不到”才是真实痛点，AI供应链控制塔可实现备件实时可见性，减少因信息缺失导致的停飞。
  \item AI驱动的“可靠性控制委员会”可识别导致大量停飞天数的长尾零件和流程，反向推动供应商增产、设计变更甚至重新审视强制检查的必要性。
\end{itemize}
\subsection{欧洲军工制造瓶颈：流程效率与供应链本地化}
\textbf{欧洲军工产能瓶颈不在需求端，而在制造端的规模化能力，需优先通过流程诊断和选择性本地化解决效率与供应链脆弱性问题。}
\begin{itemize}
  \item 欧洲军工企业多为中小企业，不愿自掏腰包研究产能，担心需求节奏变化导致资产闲置，工厂规模偏小无法简单通过加班加点扩大产出。
  \item 现有设施产能受限往往源于流程中的低效环节（如设计转量产、测试阶段），而非厂房面积，先做流程诊断再决定是否扩产可避免资本浪费。
  \item 欧洲军工从美国、亚洲采购材料和零部件，长价值链带来更长交货周期，且地缘紧张时海外供应商优先服务本国OEM，欧洲客户排后。
  \item 应对思路是“选择性本地化”，承包商和国防部需识别关键零部件并推动本土化生产，以降低供应链脆弱性。
\end{itemize}
\newpage
\section{辅助信号｜智库/国际机构}
本日IMF系列报告聚焦三大结构性主题：全球供应链脆弱性（稀土断供、关税传导、能源粮食价格传导）、人口与金融结构变迁（老龄化重塑银行、亚太金融一体化滞后）、以及财政纪律与制度能力（前沿经济体发债后困境、援助依赖型国家出路、公共投资执行率、中期预算框架）。这些研究为理解宏观风险、产业战略和资产叙事提供了非市场视角的辅助信号。
\subsection{全球供应链脆弱性与价格传导}
\textbf{从稀土断供到关税重组，再到能源粮食价格传导，IMF系列报告揭示全球供应链的结构性风险并非均匀分布，而是取决于产业网络中心度、替代弹性以及政策干预的棘轮效应。}
\begin{itemize}
  \item 稀土供应80\%中断情景下，日本GDP损失最大（-1.8\%），并非因其用量最多，而是汽车、电子等稀土密集型部门在国内生产网络中处于枢纽位置，减产通过投入产出链条广泛传导；美国损失（-1.5\%）虽低于日本，但因其汽车和电子部门的前向-后向关联更强，传导面更广（R036）。
  \item 替代弹性是决定稀土断供损失规模的单一变量，解释了低弹性和高弹性情景之间97\%的GDP损失差异；传统VAAR指标因假设零替代弹性且只算直接依赖，容易高估德国、低估美国（R036）。
  \item 2025年美国关税冲击中，出口商价格几乎未调整（品种层面零传递），关税成本完全转嫁给进口商；但产品层面平均进口价格（不含关税）反而下降，驱动力是进口商将采购从高价格国家转向低价格国家，这种“成本节约”实为质量降级（R040）。
  \item 国际油价每上涨1美元，国内汽油和柴油零售价格在3个月内传导约0.8美元；小麦传导需约10个月达到峰值0.81美元，大米传导率更低仅0.69美元；存在明显“棘轮效应”——价格上涨时传导快，下跌时传导慢，掩盖通胀累积效应（R039）。
  \item 发达经济体燃料传导接近完全，中东和北非、撒哈拉以南非洲因补贴和行政定价传导率显著偏低；2022年价格冲击中，维持低传导率的国家补贴账单在某些情况下达GDP的2\%至4\%，将调整负担转移至公共财政（R039）。
\end{itemize}
\begin{figure}[htbp]
\centering
\includegraphics[width=0.88\linewidth]{figures/fig_077.jpg}
\caption{Figure 1 - Figure 1: Eurobond Issuances by EMDEs, 2005–23 (number of issuances per year) Note: Eurobond issuances in Frontier Economies peaked in 2021, declining sharply over the next two years as global financial conditions tightened. Des}
\end{figure}
\begin{figure}[htbp]
\centering
\includegraphics[width=0.88\linewidth]{figures/fig_097.jpg}
\caption{Figure 1 - Figure 2: Distribution of within-product standard deviation of changes in U.S. tariffs over 2024-25 and 2017-19 Note: Changes in U.S. tariffs in panel (a) are computed as differences between the modal U.S. statutory tariff rat}
\end{figure}
\subsection{人口与金融结构变迁}
\textbf{亚洲老龄化正在从家庭行为端重塑银行资产负债表，而亚太金融一体化程度远落后于贸易一体化，新兴市场在全球资本流动中的“接入深度”有限，既是风险隔离也是发展瓶颈。}
\begin{itemize}
  \item 老年人口占比每上升1个百分点，银行贷款对存款比率下降约4-5个百分点；存款增长若无对应贷款需求，银行被迫将资金配置到风险更高的非贷款资产上，盈利模式和风险轮廓需重新审视（R034）。
  \item 亚太地区贡献全球三分之一贸易和GDP，但外部资产和负债占全球比重分别仅21\%和16\%；剔除中国内地与香港之间的双向头寸后，区域金融一体化真实程度更低（R041）。
  \item 亚太先进经济体（日本、澳大利亚、新加坡）金融一体化程度已接近贸易一体化水平；新兴亚太地区跨境金融资产和负债增量仅相当于GDP的6\%，为所有新兴市场区域中最低，而亚太先进地区同期为22.5\%（R041）。
  \item 亚太新兴市场更多通过贸易渠道而非资本流动渠道参与全球化；更强的金融联系会放大对全球冲击和全球金融周期的暴露，推进一体化需配套宏观审慎政策和资本流动管理工具（R041）。
\end{itemize}
\begin{figure}[htbp]
\centering
\includegraphics[width=0.88\linewidth]{figures/fig_082.jpg}
\caption{Figure 2 - Figure 3: Share of population aged 65+}
\end{figure}
\begin{figure}[htbp]
\centering
\includegraphics[width=0.88\linewidth]{figures/fig_101.jpg}
\caption{Figure 1 - Sources: IMF World Economic Outlook, CDIS, CPIS, Bank for International Settlements, and IMF staff calculations. Note: Asia-Pacific's share of global trade and cross-border financial positions are measured in terms of aggregates of imports and exports and asse}
\end{figure}
\subsection{财政纪律与制度能力}
\textbf{从前沿经济体发债后困境到援助依赖型国家的政策选择，再到公共投资执行率和中期预算框架，IMF报告一致指向制度能力而非资金规模是财政可持续性的关键变量。}
\begin{itemize}
  \item 前沿经济体首次发债后，多数国家债务状况并未改善，反而因财政纪律松弛和利息负担加重而恶化；全球流动性是“门票”，但人均收入、治理质量和外汇储备等国内基本面才是持续市场准入的“入场资格”（R032）。
  \item 援助依赖型国家面临永久性援助削减时，三种政策路径代价不同：完全依赖借款短期延续但长期加深倒退（第5年GDP低0.8\%，基尼系数从0.35升至0.45）；直接支出削减全面恶化；渐进式国内收入动员中期痛苦但长期可持续（R033）。
  \item 莱索托公共投资占GDP约10\%，远超区域平均水平，但资本支出预算执行率仅53\%，近一半项目资金闲置；问题在于项目进入预算时未经充分准备（缺乏成本测算和绩效目标），中期预算框架形同虚设（R035）。
  \item 马来西亚已通过《公共财政与财政责任法》设定赤字和债务红线，但预算仍停留在“一年一议”模式，运营与发展支出使用独立系统互不联通；IMF建议交通、卫生、教育三个部委试点建立三年期预算上限（R037）。
  \item 毛里求斯GDP核算中，全球商业公司（GBC）的交易活动尚未被纳入；IMF反对使用成本加总法，坚持应以交易价差作为产出估值基础，以保护GDP数据的国际可比性和可信度（R038）。
\end{itemize}
\begin{figure}[htbp]
\centering
\includegraphics[width=0.88\linewidth]{figures/fig_078.jpg}
\caption{Figure 2 - Figure 2: Matching Rate for Cutoff Lengths extending the access requirement beyond two decades yields little additional improvement in classification accuracy. We therefore adopt 20 years of sustained market access as the reason}
\end{figure}
\begin{figure}[htbp]
\centering
\includegraphics[width=0.88\linewidth]{figures/fig_087.jpg}
\caption{Figure 1 - Third, limited institutional capacity and weak enforcement of existing processes constrain the system's ability to correct these problems in-year, including through effective portfolio monitoring and timely reallocation. The proposed reform agenda focuses on s}
\end{figure}
\section{结语}
后续需验证的关键节点包括：中东局势对油价的实际冲击幅度及持续时间；AI资本开支货币化路径的财报信号（如云厂商资本开支回报率）；中国财政刺激落地节奏与内需修复数据；电池消费税对二线厂商盈利的实际影响；以及欧洲电网投资政策推进速度。
\newpage
\section{覆盖概览}
投行/券商 31 篇；战略咨询 3 篇；智库/国际机构 10 篇。\par
投行覆盖：GS 17篇 / Bernstein 6篇 / NOM 5篇 / JPM 2篇 / MS 1篇\par
\section{Disclaimer}
{\scriptsize\color{refgray}
This community is a paid learning and information-sharing community created for general educational purposes only. The membership fee is charged solely for access to community discussions, educational content, organization, and operational costs. It is not an advisory fee, management fee, performance fee, brokerage fee, or compensation for personalized financial, investment, legal, tax, or professional advice.\par
I participate and share information solely as an individual and for educational discussion among community members. I am not acting as a financial adviser, investment adviser, broker, dealer, portfolio manager, legal adviser, tax adviser, accountant, fiduciary, or any other licensed professional adviser. Nothing shared in this community constitutes, and should not be interpreted as, financial advice, investment advice, legal advice, tax advice, a recommendation, solicitation, offer, endorsement, instruction, or guarantee to buy, sell, hold, short, trade, or invest in any security, cryptocurrency, fund, derivative, strategy, or other financial product.\par
All content shared in this community, including messages, discussions, links, files, charts, examples, market commentary, watchlists, AI-generated or AI-polished summaries, and third-party materials, is general, impersonal, and educational in nature. It does not take into account any individual's financial situation, investment objectives, risk tolerance, investment horizon, liquidity needs, tax status, legal restrictions, or suitability requirements. No content should be relied upon as suitable for any specific person, account, portfolio, or financial decision.\par
Information shared in this community may be incomplete, inaccurate, outdated, speculative, AI-generated, AI-edited, or based on third-party internet sources that have not been independently verified. I make no representation or warranty as to the accuracy, completeness, timeliness, reliability, or suitability of any information shared.\par
Financial markets involve substantial risk, including the possible loss of principal. Past performance, historical data, hypothetical examples, simulated results, backtests, personal opinions, or educational examples do not guarantee future results. Each member is solely responsible for conducting their own research, exercising independent judgment, and making their own decisions. Before making any financial, investment, legal, or tax decision, members should consult qualified and licensed professionals.\par
Participation in this community, payment of any membership fee, or communication with me or other members does not create any adviser-client, fiduciary, professional, contractual, agency, partnership, employment, or other advisory relationship. By joining or participating in this community, you acknowledge and agree that you are solely responsible for your own decisions, actions, transactions, profits, losses, risks, and consequences, and that I shall not be liable for any direct, indirect, incidental, consequential, financial, legal, tax, or other loss or damage arising from or related to any content, discussion, information, or material shared in this community.\par
}
\newpage
\begin{center}
{\Large 更多详情报告kcdesk.com}\par\vspace{0.8cm}
\includegraphics[width=0.58\linewidth]{\detokenize{../../prompts/zsxq_img.jpg}}
\end{center}
\end{document}
